\section{Setup and notation}\label{sec-setup}

\begin{objectif}
Fix the notation and establish the three tools the rest of the note needs: the forward formula, the carry
difference, and the passage from Totem data to the implied repo.
\end{objectif}

\subsection{Conventions}
Time is measured in years from the valuation date, and rates are continuously compounded zero rates. Scalars and
vectors are written in lower case, matrices in upper case; $\tr{A}$ is the transpose of $A$, $\M_{N,m}(\R)$ the space
of matrices with $N$ rows and $m$ columns, $\sprod{x}{y}=\tr{x}\,y$ the inner product of $\R^N$ and $\norm{x}$ the
associated norm. The superscript $\tot$ marks a quantity from Totem and the superscript $\fo$ a front-office quantity.
A basis point (bp) is $10^{-4}$; a carry difference is expressed in basis points of spot. The notation is summarised
in \cref{app-notation}.

\subsection{Forward and discrete dividends}
Let $S_0$ be the spot, $r(T)$ the zero rate of maturity $T$, and $q(T)$ the zero \emph{repo} rate, that is, the cost
of borrowing the stock: holding the stock financed by a repo earns $r-q$ per unit of time. Dividends are discrete: an
amount $d_i$ goes ex at time $t_i$. Let $\tilde d_i=d_i\,e^{-(r(t_i)-q(t_i))\,t_i}$ be the dividend discounted at the
carry rate.

\begin{definition}[Forward]\label{def-forward}
The forward of maturity $T$ is
\begin{equation}\label{eq-forward}
F(T)=\bigl(S_0-D(T)\bigr)\,e^{(r(T)-q(T))\,T}, \qquad D(T)=\sum_{t_i\le T}\tilde d_i .
\end{equation}
$D(T)$ is the sum of the dividends going ex before $T$, discounted at the carry rate.
\end{definition}

This is the usual replication argument: buy the stock financed at the rate $r-q$ and carry the dividends received
forward. The repo also enters $D$, weakly, through the discounting of the dividends; the bootstrap of
\cref{sec-bootstrap} handles this coupling exactly.

\begin{lemma}[Log of the forward]\label{lem-logf}
$\ln F(T)=\bigl(r(T)-q(T)\bigr)\,T+\ln\bigl(S_0-D(T)\bigr)$. Hence, to first order, perturbations of the forward, the
dividends and the repo satisfy
\begin{equation}\label{eq-additivity}
\delta\bigl(q(T)\,T\bigr)=\delta\bigl(r(T)\,T\bigr)-\delta\ln F(T)-\frac{\delta D(T)}{S_0-D(T)} .
\end{equation}
\end{lemma}
\begin{proof}
Take the log of \eqref{eq-forward} and differentiate, with $\delta\ln(S_0-D)=-\delta D/(S_0-D)+O(\delta D^2)$.
\end{proof}

\subsection{The carry difference}
The Totem maturities, or \emph{pillars}, are $T_1<\dots<T_N$ (here $N=11$, from one month to five years). For any
quantity $x$, we write $x_k=x(T_k)$.

\begin{definition}[Carry difference]\label{def-delta}
The \emph{carry difference} at pillar $T_k$ is
\begin{equation}\label{eq-delta}
\delta_k=\bigl(q^{\tot}_k-q^{\fo}_k\bigr)\,T_k, \qquad k\in\ient1N,
\end{equation}
where $q^{\tot}$ is the Totem-implied repo (\cref{def-implied}) and $q^{\fo}$ the FO repo. We write
$\delta=(\delta_1, \dots, \delta_N)\in\R^N$.
\end{definition}

\begin{proposition}[The repo, a difference of forward and dividends]\label{prop-delta}
The controller uses the same rate curve on both sides. Then, up to the coupling through the discounting of the
dividends,
\begin{equation}\label{eq-delta-exact}
\delta_k=-\ln\frac{F^{\tot}_k}{F^{\fo}_k}+\ln\frac{S_0-D^{\tot}_k}{S_0-D^{\fo}_k}
\;\approx\;-\frac{F^{\tot}_k-F^{\fo}_k}{F_k}-\frac{D^{\tot}_k-D^{\fo}_k}{S_0-D_k} .
\end{equation}
\end{proposition}
\begin{proof}
By construction, the Totem-implied repo reprices the Totem forward with the Totem dividends, and the FO repo the FO
forward with the FO dividends. Apply \cref{lem-logf} on both sides and subtract: the rate terms cancel. The
approximation comes from $\ln(1+x)=x+O(x^2)$.
\end{proof}

In carry terms, the repo difference is therefore the relative forward difference minus the relative dividend
difference. This is already, in elementary form, the controller's intuition: repo control $=$ forward control $-$
dividend control. \Cref{sec-control} turns it into an identity on the control lines.

\begin{retenir}
Why carry rather than repo: in $\delta$, the effects of the different mechanisms \emph{add up}, and nothing is
magnified. A constant error $\varepsilon$ on $\ln F$ is a constant in carry and a hyperbola $\varepsilon/T$ in repo.
Every test in this note therefore works on $\delta$.
\end{retenir}

\subsection{The Totem data}
\begin{assumption}[Consistent contributors]\label{hyp-coherent}
Contributor $j$ has a repo view $q_j(\cdot)$ and a dividend view $D_j(\cdot)$, and quotes option prices consistent with
the forward $F_j(T)=\bigl(S_0-D_j(T)\bigr)\,e^{(r(T)-q_j(T))T}$, the spot and the rate curve being common.
\end{assumption}

\begin{assumption}[Panels]\label{hyp-panels}
At pillar $T_k$, the option panel $\Popt_k$ has $n^{\opt}_k$ contributors; for period $p$, the dividend panel $\Pdiv_p$
has $n^{\dv}_p$. Panels are nested in maturity: a contributor who quotes a long maturity also quotes the shorter ones.
View errors are independent from one contributor to another. The two services have $n^{\com}$ contributors in common
($n^{\com}=0$ for disjoint panels).
\end{assumption}

\begin{assumption}[Consensus]\label{hyp-consensus}
The consensus is the arithmetic mean over the panel (after an optional outlier filter). For each maturity and each
period, Totem publishes the consensus, the dispersion across contributors and their number.
\end{assumption}

\subsection{From the consensus to the implied repo}\label{sec-bootstrap}
The passage has three steps.

\paragraph{Forwards from put--call parity.}
\begin{lemma}[Parity and consensus forward]\label{lem-parity}
If every contributor quotes European calls and puts with the same discount factor $B(0, T)$, then
$C_j(K)-P_j(K)=B(0, T)\,(F_j-K)$ for every strike $K$. If the consensus averages the same contributors at every
strike, $\bar C(K)-\bar P(K)=B(0, T)\,(\bar F-K)$, where $\bar F$ is the average of the contributors' forwards. The
strike where $\bar C=\bar P$ is therefore $\bar F$, and it does not depend on $B(0, T)$.
\end{lemma}
\begin{proof}
The first equality is put--call parity; the second follows by averaging the first, which is affine; the third by
setting the second to zero.
\end{proof}
In practice, the line $K\mapsto\bar C(K)-\bar P(K)$ is fitted by weighted least squares; its root is the Totem forward
$F^{\tot}_k$ and its slope gives an implied discount factor, used as a diagnostic. Two situations break the lemma: an
outlier filter applied separately per strike and option type (the averaged contributors are then no longer the same),
and American exercise. A worked example and the case of discount factors that differ across contributors are in
\cref{app-parity}.

\paragraph{Dividend allocation.} The Totem levels $\ell_p$ are sums of the dividends going ex in the period
$I_p=\intoc{e_{p-1}}{e_p}$, $p\in\ient1m$ (here $m=6$: the rest of 2026, then 2027 to 2031). The controller spreads them
over the ex-dates of the FO forecast: a profile gives a weight $\pi_i\ge0$ to each date, with
$\sum_{t_i\in I_p}\pi_i=1$, and the allocated dividend is $d_i=\ell_p\,\pi_i$.

\begin{lemma}[Allocation]\label{lem-allocation}
For a fixed repo,
\[
D_k=\sum_{p=1}^m a_{kp}\,\ell_p, \qquad a_{kp}=\sum_{t_i\in I_p,\ t_i\le T_k}\pi_i\,e^{-(r(t_i)-q(t_i))\,t_i} .
\]
We call $A=(a_{kp})\in\M_{N,m}(\R)$ the allocation matrix.
\end{lemma}
\begin{proof}
Replace $d_i$ by $\ell_p\pi_i$ in \eqref{eq-forward} and group the terms by period.
\end{proof}
The coefficient $a_{kp}$ is the discounted share of period $p$ that goes ex before $T_k$. It is zero if the period
starts after $T_k$, and equal to the discount factor of the period (close to $1$) if the whole period goes ex before
$T_k$; in between, for a period that straddles a pillar, it depends on the profile: this is exactly where the profile
matters.

\paragraph{Bootstrap.} We look for a repo curve with a constant forward repo between pillars, that is, such that
$t\mapsto q(t)\,t$ is affine on each $[T_{k-1}, T_k]$, which reprices the forwards. It is built pillar by pillar:
knowing $q_1, \dots, q_{k-1}$, the unknown is $x=q_kT_k$, and the model forward at $T_k$ is a function
$\mathcal F_k(x)$, written out in \cref{app-bootstrap}.

\begin{proposition}[The bootstrap is well posed]\label{prop-bootstrap}
Assume $S_0-D_{k-1}>0$. The function $\mathcal F_k$ satisfies $\mathcal F_k'\le-\mathcal F_k$; it is therefore
strictly decreasing where it is positive. For every $F_k>0$, the equation $\mathcal F_k(x)=F_k$ has a unique solution,
and Newton's method applied to $x\mapsto\ln\mathcal F_k(x)-\ln F_k$, safeguarded by bisection, converges to it.
\end{proposition}
The proof, a worked example and the algorithm are in \cref{app-bootstrap}.

\begin{definition}[Totem-implied repo]\label{def-implied}
The \emph{Totem-implied repo} $q^{\tot}$ is the bootstrap solution for the forwards $F^{\tot}_k$ and the Totem
dividend levels $\ell^{\tot}_p$ allocated on the FO profile. We also write $\hat q$ for the solution with the Totem
forwards and the \emph{FO} dividends: it depends on Totem only through its forwards.
\end{definition}

\begin{retenir}
\begin{itemize}[leftmargin=12pt]
\item $\delta_k=(q^{\tot}_k-q^{\fo}_k)\,T_k$: the repo difference expressed in carry, in basis points of spot.
\item $\delta_k\approx-$ (relative forward difference) $-$ (relative dividend difference): the implied repo is a
residual.
\item The passage from Totem to the implied repo (parity, allocation, bootstrap) is exact and well posed; the
difficulties come from the data, not from the algorithm.
\end{itemize}
\end{retenir}
